% !TEX root = ../../ctfp-print.tex

\lettrine[lhang=0.17]{程}{序员们围绕单子(monad)}发展出了一整套神话体系。它被认为是编程领域最抽象最难懂的概念之一。有人「顿悟」了单子的本质，也有人始终不得其门而入。对许多人来说，理解单子概念的瞬间就像经历了一场神秘主义体验。单子抽象了如此多样的构造本质，以至于我们在日常生活中根本找不到恰当的类比。我们只能像盲人摸象般在黑暗中摸索：有人喊着「这是根绳子」，有人高呼「这是树干」，还有人大叫「这是个卷饼！」

让我来澄清事实：关于单子的所有神秘感都源于误解。单子本质上是个非常简单的概念。造成困惑的根源恰恰在于它的应用多样性。

为了撰写本文，我特意查阅了电工胶带（又称鸭绒胶带）的应用场景。以下是它能完成的一些典型任务：

\begin{itemize}
  \tightlist
  \item
        密封管道
  \item
        修复阿波罗13号上的二氧化碳洗涤器
  \item
        治疗疣
  \item
        解决苹果iPhone 4掉话问题
  \item
        制作毕业舞会礼服
  \item
        搭建悬索桥
\end{itemize}

\noindent
现在假设你不知道什么是电工胶带，仅凭上述清单去理解这个概念——祝你好运！

因此我想给「单子就像...」的陈词滥调再添一笔：单子就像电工胶带。它的应用场景千差万别，但核心原理极其简单：它能把事物粘合在一起。更确切地说，是将事物组合起来。

这在一定程度上解释了为何很多程序员（特别是从命令式语言背景转来的）难以理解单子。问题在于，我们不习惯用函数组合的方式思考编程。这种思维模式并不常见，我们习惯给中间值命名而不是直接传递结果，倾向于内联简短的粘合代码而非抽象成辅助函数。看这个C语言实现的向量长度函数：

\begin{snip}{cpp}
double vlen(double * v) {
    double d = 0.0;
    int n;
    for (n = 0; n < 3; ++n)
        d += v[n] * v[n];
    return sqrt(d);
}
\end{snip}
对比一下Haskell中显式体现函数组合的版本：

\src{snippet01}
（这里为了让事情更晦涩些，我通过将第二个参数设为\code{2}，部分应用了指数运算符\code{(\^{})}。）

我并不是说Haskell的无点(point-free)风格总是更优，只是想强调函数组合是我们编程行为的底层逻辑。虽然实际上我们也在进行函数组合，但Haskell还是提供了名为\code{do}记法的命令式语法糖来支持单子组合。稍后我们会看到具体应用。不过首先，让我解释为什么我们需要单子组合。

\section{克莱斯利范畴}

我们之前通过对常规函数进行装饰得到了
\hyperref[kleisli-categories]{写者单子(writer monad)}。这种特定的装饰方式是通过将返回值与字符串配对实现的，更一般地，是与某个幺半群(monoid)的元素配对。现在我们可以认识到，这种装饰本质上是一个函子(functor)：

\src{snippet02}
随后我们找到了一种组合装饰函数的方法，即克莱斯利箭头(Kleisli arrows)，这些函数具有如下形式：

\src{snippet03}
日志的累积正是在这样的组合过程中实现的。

现在我们可以给出克莱斯利范畴的更一般定义。给定一个范畴$\cat{C}$和一个自函子(endofunctor)$m$，相应的克莱斯利范畴$\cat{K}$与$\cat{C}$有相同对象，但态射(morphism)不同。在$\cat{K}$中，对象$a$到$b$的态射在原始范畴$\cat{C}$中表现为：
$$a \to m\ b$$
需要特别注意的是，我们把$\cat{K}$中的克莱斯利箭头视为$a$到$b$的态射，而不是$a$到$m\ b$的态射。

在之前的例子中，$m$特化成了\code{Writer w}，其中\code{w}是某个固定的幺半群。

只有当能够定义合适的组合方式时，克莱斯利箭头才能构成范畴。如果存在满足结合律且每个对象都有单位箭头的组合操作，那么这个函子$m$就被称为\newterm{单子(monad)}，对应的范畴称为克莱斯利范畴。

在Haskell中，克莱斯利组合通过鱼算符(fish operator)\code{>=>}定义，单位箭头是多态函数\code{return}。下面是基于克莱斯利组合的单子定义：

\src{snippet04}
需要注意的是，单子有多种等价定义方式，这并非Haskell生态中最主流的形式。我喜欢这个定义是因为其概念简洁且直观清晰，但在实际编程中可能有更便捷的定义方式。我们马上就会讨论这些替代方案。

在这种表述下，单子定律(monad laws)非常容易表达。虽然Haskell无法强制实施这些定律，但它们可用于等式推理。这些定律其实就是克莱斯利范畴的标准组合定律：

\begin{snip}{haskell}
(f >=> g) >=> h = f >=> (g >=> h) -- 结合律
return >=> f = f                  -- 左单位元
f >=> return = f                  -- 右单位元
\end{snip}
这种定义方式完美诠释了单子的本质：它是一种组合装饰函数的方式。这不是关于副作用或状态的机制，而是关于组合本身的抽象。正如我们稍后将看到的，装饰函数可以用来表达各种效应或状态，但这并非单子的核心功能。单子就是那条黏性的电工胶带，把一个装饰函数的输出端牢牢粘住另一个装饰函数的输入端。

回到我们的\code{Writer}示例：只要\code{w}满足幺半群定律（同样无法在Haskell中强制实施），\code{Writer}的日志记录函数（即克莱斯利箭头）就能构成范畴，因为此时\code{Writer}已经是一个单子：

\src{snippet05}

\code{Writer}单子还有一个很有用的克莱斯利箭头叫做\code{tell}，它的唯一作用就是将参数添加到日志中：

\src{snippet06}
我们将把这个基本构件用于构建其他单子函数。

\section{鱼操作符的解剖}

在为不同单子（monad）实现鱼操作符时，你会迅速发现大量代码重复且可以轻松提取复用。
首先，两个函数的克莱斯利复合（Kleisli composition）必须返回一个函数，因此其实现可以从接收类型为\code{a}参数的lambda表达式开始：

\src{snippet07}
我们唯一能做的就是将此参数传递给\code{f}：

\src{snippet08}
此时我们需要产生类型为\code{m c}的结果，而我们拥有类型为\code{m b}的对象和函数\code{g :: b -> m c}。让我们定义一个完成此功能的函数，该函数被称为\emph{绑定（bind）}，通常以中缀操作符形式书写：

\src{snippet09}
对于每个单子，我们可以选择定义bind而非鱼操作符。实际上标准Haskell单子定义正是使用bind：

\src{snippet10}
以下是\code{Writer}单子的bind定义：

\src{snippet11}
确实比鱼操作符的定义更简洁。

我们可以通过分解bind进一步探索，利用\code{m}作为函子（functor）的特性。我们可以使用\code{fmap}将函数\code{a -> m b}应用于\code{m a}的内容，这会将\code{a}转换为\code{m b}。因此应用后的结果类型为\code{m (m b)}。虽然这不是我们想要的\code{m b}类型，但我们已经很接近了。我们只需要一个能折叠或展平双重\code{m}应用的函数，这样的函数称为\code{join}：

\src{snippet12}
通过\code{join}，我们可以将bind重写为：

\src{snippet13}
这引导我们得到定义单子的第三种方式：

\src{snippet14}
此处我们显式要求\code{m}必须是\code{Functor}。在前两种单子定义中无需这样做，因为任何支持鱼操作符或bind操作符的类型构造器\code{m}自动就是函子。例如，可以用bind和\code{return}定义\code{fmap}：

\src{snippet15}
为了完整性，以下是\code{Writer}单子的\code{join}实现：

\src{snippet16}

\section{do表示法}

使用单子（monad）编写代码的一种方式是通过Kleisli箭头---使用fish算子进行组合。这种编程模式是对无点（point-free）风格的推广。无点代码紧凑且通常非常优雅，但总体而言可能难以理解，近乎晦涩。这就是大多数程序员倾向于为函数参数和中间值命名的原因。

当处理单子时，这意味着优先选择bind算子而非fish算子。bind接受一个单子值并返回一个单子值。程序员可以选择为这些值命名，但这几乎没有改进。我们真正想要的是假装自己在处理常规值，而不是封装它们的单子容器。这正是命令式代码的工作方式---诸如更新全局日志之类的副作用大多被隐藏起来。而这就是Haskell中\code{do}表示法模拟的效果。

您可能会疑惑，既然如此为何还要使用单子？如果我们想让副作用不可见，为什么不直接使用命令式语言？答案在于单子让我们对副作用有了更好的控制。例如，\code{Writer}单子中的日志在函数间传递而从未暴露为全局变量。这消除了日志混乱或数据竞争的可能性。此外，单子代码与程序其他部分有清晰的界限划分。

\code{do}表示法本质上只是单子组合的语法糖。表面上它看起来很像命令式代码，但实际上会转换为一系列bind操作和lambda表达式。

例如，考虑我们之前用以展示\code{Writer}单子中Kleisli箭头组合的例子。使用当前定义，它可以重写为：

\src{snippet17}
此函数将输入字符串的所有字符转为大写并拆分为单词，同时产生操作日志。

在\code{do}表示法中则呈现为：

\src{snippet18}
此处\code{upStr}只是一个\code{String}，尽管\code{upCase}产生的是\code{Writer}：

\src{snippet19}
这是因为编译器会将\code{do}块去糖化为：

\src{snippet20}
\code{upCase}的单子结果绑定到接收\code{String}的lambda表达式。正是这个字符串的名字出现在\code{do}块中。当我们读取：

\src{snippet21}
我们说\code{upStr} \emph{获取}\code{upCase s}的结果。

当我们内联\code{toWords}时，这种伪命令式风格更加明显。我们将其替换为调用记录字符串\code{"toWords "}的\code{tell}，接着调用携带分割字符串\code{upStr}结果的\code{return}。请注意\code{words}是一个处理字符串的常规函数。

\src{snippet22}
在此处，do块的每行在去糖化代码中都引入了一个新的嵌套bind：

\src{snippet23}
注意到\code{tell}产生单位值，因此无需传递给后续lambda。忽略单子结果的具体内容（但保留其效应---此处指日志贡献）是很常见的做法，因此存在专门的操作符来替代这种情况下的bind：

\src{snippet24}
我们的代码实际去糖化效果如下：

\src{snippet25}
一般而言，do块由若干行（或子块）组成：要么使用左箭头引入新名称供后续代码使用，要么纯粹出于副作用执行。bind操作符隐式存在于各行之间。顺便提及，在Haskell中可以用花括号和分号替代do块的缩进格式。这为"将单子视为重载分号机制"的说法提供了依据。

值得注意的是，当去糖化do表示法时，lambda和bind操作符的嵌套结构会产生这样的效果：后续do块的执行会基于每行的结果。这一特性可用于引入复杂控制结构，例如模拟异常处理。

有趣的是，do表示法的等价物已在C++等命令式语言中找到应用。这里指的是可恢复函数（resumable functions）或协程（coroutines）。众所周知C++
\urlref{https://bartoszmilewski.com/2014/02/26/c17-i-see-a-monad-in-your-future/}{future对象构成单子}，
这是延续体（continuation）单子的一个实例，稍后我们将讨论。延续体的问题在于其组合难度极高。在Haskell中，我们使用do表示法将"我的回调触发你的回调"这类混乱代码转化为看似顺序执行的形式。可恢复函数使C++也能实现相同转换。同样的机制还可用于将
\urlref{https://bartoszmilewski.com/2014/04/21/getting-lazy-with-c/}{嵌套循环的混乱结构}
转化为列表推导式或"生成器"，这本质就是列表单子的do表示法。如果没有统一的单子抽象，这些问题通常需要各自提供定制化的语言扩展来解决。而在Haskell中，这一切都可以通过库来实现。